% first_to_call_fifty.tex  (fixed: removed 'enhanced' key from tcolorbox styles)
\documentclass[11pt]{article}
\usepackage[margin=1in]{geometry}
\usepackage{amsmath,amssymb,amsthm}
\usepackage{booktabs}
\usepackage{microtype}
\usepackage{parskip}
\usepackage{tcolorbox}
\tcbuselibrary{breakable}
\usepackage{hyperref}
\hypersetup{colorlinks=true, linkcolor=blue!60!black, urlcolor=blue!60!black}

\title{\bfseries First to Call 50\\[4pt] \large Winning strategy using backward induction}
\author{Souvik Sarkar \\ \texttt{sounix000@gmail.com}}
\date{April 13, 2026}

% tcolorbox presets for consistent boxes
% NOTE: do not use the 'enhanced' key to remain compatible with older tcolorbox versions.
\tcbset{
  mybox/.style={
    breakable,
    colback=blue!3!white,
    colframe=blue!60!black,
    boxrule=0.7pt,
    left=6pt, right=6pt, top=6pt, bottom=6pt,
    fonttitle=\bfseries,
    arc=3pt
  },
  resultbox/.style={
    breakable,
    colback=green!5!white,
    colframe=green!50!black,
    boxrule=0.8pt,
    left=6pt, right=6pt, top=6pt, bottom=6pt,
    fonttitle=\bfseries,
    arc=3pt
  },
  conclbox/.style={
    breakable,
    colback=gray!5!white,
    colframe=black!40!gray,
    boxrule=0.6pt,
    left=6pt, right=6pt, top=6pt, bottom=6pt,
    fonttitle=\bfseries,
    arc=3pt
  }
}

\begin{document}
\maketitle

\section*{Problem statement}

\begin{tcolorbox}[mybox, title=Game setup]
Two players take turns calling out integers. The first person to call out \textbf{50} wins.

\medskip
\textbf{Rules:}
\begin{itemize}
  \item The first player must call an integer between $1$ and $10$ (inclusive).
  \item Each subsequent number must be strictly greater than the previous number by at least $1$ and at most $10$.
  \item The player who first calls exactly $50$ wins immediately.
\end{itemize}

\noindent\textbf{Question:} Should you go first, and what is the optimal strategy?
\end{tcolorbox}

\section*{Informal analysis}

Let the calls during the game be denoted
\[
C_1, C_2, \dots, C_w,
\]
where $C_1$ is the first call and $C_w = 50$ is the winning call. The rules imply
\[
1 \le C_1 \le 10,\qquad
1 \le C_i - C_{i-1} \le 10 \quad\text{for } i\ge 2.
\]

We analyze the game by backward induction: identify positions (numbers) that are \emph{winning positions} for the player about to move, meaning that from such a number the player to move can force a win with perfect play.

\section*{Backward induction (key observation)}

To call $50$ on your turn, the previous call must lie in the range $[40,49]$. Therefore, if you can force the opponent to hand you any number in $[40,49]$, you will win by calling $50$.

To force the opponent into $[40,49]$, you would like to hand them the number $39$, because from $39$ the opponent's allowed responses are $[40,49]$ and then you can call $50$.

Repeating this backward:
\[
39,\; 28,\; 17,\; 6
\]
are the critical numbers: if you call one of these on your turn, the opponent is forced into an interval that allows you to move to the next critical number (adding $11$ each time), eventually reaching $39$ and then $50$.

Thus the numbers
\[
6,\, 17,\, 28,\, 39,\, 50
\]
form the thread of winning positions spaced $11$ apart.

\section*{Formal justification}

\subsection*{Claim}
If a player on their turn calls a number congruent to $6\pmod{11}$ (i.e., of the form $6+11k$), then with perfect play that player can force a win; conversely, if the player to move holds a number not congruent to $6\pmod{11}$, the opponent can (by proper play) move to a number congruent to $6\pmod{11}$.

\subsection*{Proof}
Suppose on your turn you call $n = 6 + 11k$ for some integer $k\ge 0$ (so $n\le 50$ as relevant). The opponent must respond with a number $m$ satisfying
\[
n+1 \le m \le n+10.
\]
All such $m$ satisfy
\[
m \equiv n + r \pmod{11} \quad\text{for some } r\in\{1,\dots,10\},
\]
so $m \not\equiv 6 \pmod{11}$. In particular, no response of the opponent can be congruent to $6\pmod{11}$. Therefore, you can always respond by calling
\[
n' = n + 11 = 6 + 11(k+1),
\]
which is legal because $n' - m \le n' - (n+1) = 10$ and $n' - m \ge n' - (n+10) = 1$. Thus you move to the next number in the sequence $6,17,28,39,50$.

This induction continues until you call $50$ (when $n'$ reaches $50$). Hence control of numbers congruent to $6\pmod{11}$ lets you force the win.

Conversely, if it is your opponent's turn and they hold a number not congruent to $6\pmod{11}$, you can always choose a response (within $1$--$10$ added) that lands on the next $6\pmod{11}$ number, so the other player can be forced back into the losing class.

Therefore the positions congruent to $6\pmod{11}$ are precisely the winning positions for the player about to move.

\section*{Concrete strategy}

\begin{tcolorbox}[resultbox, title=Optimal strategy]
Go first and call \(\mathbf{6}\). After that, whenever your opponent calls some number, respond by playing the unique number that brings the running total to the next term of the sequence
\[
6,\;17,\;28,\;39,\;50,
\]
i.e., always move to the smallest number of the form \(6+11k\) that is at most \(50\). This guarantees that you will be the one to call \(50\).
\end{tcolorbox}

\section*{Winning sequence (example)}

The following table shows the moves if both players follow the described structure (opponent's calls are shown as allowed ranges):

\begin{center}
\renewcommand{\arraystretch}{1.15}
\begin{tabular}{@{}lllll@{}}
\toprule
Turn & Player & My call & Opponent's possible responses & Notes \\
\midrule
1 & Me       & 6  & ---          & start at a \(6\pmod{11}\) position \\
2 & Opponent & ---& $[7,16]$    & any call in this interval \\
3 & Me       & 17 & ---          & move to next \(6\pmod{11}\) number \\
4 & Opponent & ---& $[18,27]$   & \\
5 & Me       & 28 & ---          & \\
6 & Opponent & ---& $[29,38]$   & \\
7 & Me       & 39 & ---          & \\
8 & Opponent & ---& $[40,49]$   & \\
9 & Me       & 50 & ---          & win \\
\bottomrule
\end{tabular}
\end{center}

\section*{Why the strategy works}

Two simple facts suffice:
\begin{enumerate}
  \item On any move you can increase the current number by $1$--$10$.
  \item The residue classes modulo $11$ partition the integers into blocks of size $11$. By forcing the running total into the residue class $6\pmod{11}$ on your turns, you deny the opponent the ability to do so on theirs; therefore you can maintain control and eventually reach $50$.
\end{enumerate}

\section*{Conclusion}

\begin{tcolorbox}[conclbox, title=Summary]
The first player wins with perfect play by starting at \(\mathbf{6}\) and then repeatedly adding \(11\) (after the opponent moves), following the sequence
\[
6,\;17,\;28,\;39,\;50.
\]
Thus the winning positions are exactly those congruent to \(6\pmod{11}\).
\end{tcolorbox}

\vspace{6pt}
\noindent\textit{Remark.} The same backward-induction method generalizes: if on each move a player may add between \(1\) and \(M\), and the target is \(T\), then winning positions are those congruent to \(r\pmod{M+1}\) where \(r\) is the smallest positive integer such that \(T\equiv r\pmod{M+1}\) and that is reachable from a legal initial move.

\end{document}
